\documentclass[a4paper,12pt]{article}

% ===========================================================================
%  Лабораторная работа 3.7.3 «Длинные линии»
%
%  Сборка:  cd 3.7.3/report && pdflatex main.tex  (дважды, или latexmk -pdf)
%  Таблицы, графики и числа читаются из ../data/processed/,
%  которые готовит ../scripts/process_data.py.
%
%  До измерений отчёт собирается: таблицы пустые, на месте чисел -- «???»,
%  на месте осциллограмм -- рамки-заглушки с именем ожидаемого файла.
% ===========================================================================

% ============================================================
% Общий преамбул для отчётов по лабораторным работам (labs_lpr)
% Подключается из report/main.tex командой:
%   % ============================================================
% Общий преамбул для отчётов по лабораторным работам (labs_lpr)
% Подключается из report/main.tex командой:
%   % ============================================================
% Общий преамбул для отчётов по лабораторным работам (labs_lpr)
% Подключается из report/main.tex командой:
%   \input{../../common/latex_template/preamble.tex}
% ============================================================

% Поддержка русского языка и кодировки
\usepackage[T2A]{fontenc}
\usepackage{lmodern}
\usepackage[utf8]{inputenc}
\usepackage[russian]{babel}

% Математические пакеты
\usepackage{amsmath, amssymb}

% Настройка полей
\usepackage{geometry}
\geometry{top=2cm, bottom=2cm, left=2.5cm, right=1.5cm}

% Колонтитулы
\usepackage{fancyhdr}
\usepackage{graphicx}

% Современные таблицы
\usepackage{tabularray}

% Работа с csv файлами
\usepackage{csvsimple-l3}
\usepackage{tikz}

% Графики
\usepackage{pgfplots}
\pgfplotsset{compat=1.18}

% Ссылки
% \usepackage{hyperref} % гипперссылки, нужно оформление вместо красного квадрата
\usepackage[russian]{cleveref} % умные ссылки, нужно переоформление для формул, рисунков и тд (вместо ф-л что-то нужно придумать)

% ============================================================

% Поддержка русского языка и кодировки
\usepackage[T2A]{fontenc}
\usepackage{lmodern}
\usepackage[utf8]{inputenc}
\usepackage[russian]{babel}

% Математические пакеты
\usepackage{amsmath, amssymb}

% Настройка полей
\usepackage{geometry}
\geometry{top=2cm, bottom=2cm, left=2.5cm, right=1.5cm}

% Колонтитулы
\usepackage{fancyhdr}
\usepackage{graphicx}

% Современные таблицы
\usepackage{tabularray}

% Работа с csv файлами
\usepackage{csvsimple-l3}
\usepackage{tikz}

% Графики
\usepackage{pgfplots}
\pgfplotsset{compat=1.18}

% Ссылки
% \usepackage{hyperref} % гипперссылки, нужно оформление вместо красного квадрата
\usepackage[russian]{cleveref} % умные ссылки, нужно переоформление для формул, рисунков и тд (вместо ф-л что-то нужно придумать)

% ============================================================

% Поддержка русского языка и кодировки
\usepackage[T2A]{fontenc}
\usepackage{lmodern}
\usepackage[utf8]{inputenc}
\usepackage[russian]{babel}

% Математические пакеты
\usepackage{amsmath, amssymb}

% Настройка полей
\usepackage{geometry}
\geometry{top=2cm, bottom=2cm, left=2.5cm, right=1.5cm}

% Колонтитулы
\usepackage{fancyhdr}
\usepackage{graphicx}

% Современные таблицы
\usepackage{tabularray}

% Работа с csv файлами
\usepackage{csvsimple-l3}
\usepackage{tikz}

% Графики
\usepackage{pgfplots}
\pgfplotsset{compat=1.18}

% Ссылки
% \usepackage{hyperref} % гипперссылки, нужно оформление вместо красного квадрата
\usepackage[russian]{cleveref} % умные ссылки, нужно переоформление для формул, рисунков и тд (вместо ф-л что-то нужно придумать)


% Пакеты, нужные именно этой работе
\usepackage{subcaption}
\usepackage{circuitikz}

% --- Служебные макросы ------------------------------------------------------
% Короткий пропуск (число/слово, которое надо вписать после измерений)
\newcommand{\qq}{\textbf{???}}
% Пропуск, требующий осмысленного текста
\newcommand{\fillme}{\textbf{?? НЕОБХОДИМО ЗАПОЛНИТЬ ??}}
% Имя файла моноширинным шрифтом
\newcommand{\filename}[1]{\texttt{\detokenize{#1}}}
% Скриншот из report/img/; если файла ещё нет -- рамка-заглушка.
\newcommand{\labimage}[2]{%
    \IfFileExists{img/#1}{%
        \includegraphics[width=#2\linewidth]{img/#1}%
    }{%
        \fbox{\parbox[c][0.6\dimexpr#2\linewidth\relax][c]{%
                \dimexpr#2\linewidth-2\fboxsep-2\fboxrule\relax}{%
            \centering\small\filename{img/#1}\\[2mm]\fillme}}%
    }%
}
% Итоговые числа из ../data/processed/results.tex: $\res{Vph}$.
% Если обработка ещё не запускалась -- «???».
\InputIfFileExists{../data/processed/results.tex}{}{}
\makeatletter
\newcommand{\res}[1]{%
    \ifcsname res@#1\endcsname\csname res@#1\endcsname\else\qq\fi}
% То же число без оформления -- для формул pgfplots (0, если нет данных)
\newcommand{\resnum}[1]{%
    \ifcsname resnum@#1\endcsname\csname resnum@#1\endcsname\else 0\fi}
\makeatother

% Колонтитулы (как в остальных работах репозитория)
\pagestyle{fancy}
\fancyhf{}
\fancyhead[L]{%
    \IfFileExists{../../common/latex_template/miptlogo.png}{%
        \includegraphics[height=1cm]{../../common/latex_template/miptlogo.png}%
    }{%
        \textbf{[ЛОГОТИП МФТИ]}%
    }%
}
\fancyhead[R]{\textbf{Емельянов ДА}}
\fancyfoot[C]{\thepage}
\renewcommand{\headrulewidth}{0.4pt}
\setlength{\headheight}{33pt}

% Общий стиль графиков
\pgfplotsset{
    labplot/.style={
        width=0.8\textwidth,
        height=0.5\textwidth,
        grid=both,
        minor tick num=1,
        legend cell align=left,
        legend style={font=\small},
        scaled ticks=false,
        unbounded coords=jump,
        tick label style={/pgf/number format/use comma},
    },
}

\begin{document}

% --- ТИТУЛЬНЫЙ ЛИСТ ---
\begin{titlepage}
\begin{center}
    {\large МОСКОВСКИЙ ФИЗИКО-ТЕХНИЧЕСКИЙ ИНСТИТУТ}
\end{center}
\begin{center}
    {\large Физтех-школа физики и исследований им. Ландау\break ЛФИ}
\end{center}

    \vspace{3cm}

\vspace{0.15cm}
{
\begin{center}
    \huge \bf Лабораторная работа №3.7.3 по курсу Общая физика\\
    \vspace{0.4cm}
    \LARGE Длинные линии.
\end{center}
}
\vspace{1.5cm}
\begin{flushright}
Емельянов Дмитрий Алексеевич\\
\vspace{0.4cm}
Б02-502
\end{flushright}
\vspace{10cm}
\begin{center}
    \large 06 октября 2026, г. Долгопрудный
\end{center}
\end{titlepage}
% ----------------------

\section{Аннотация}

\textbf{Цель работы:} ознакомиться и проверить на практике теорию
распространения электрических сигналов вдоль длинной линии; измерить
амплитудно- и фазово-частотные характеристики коаксиальной линии; определить
погонные характеристики такой линии; на примере модели длинной линии изучить
распределение амплитуды колебаний сигнала по длине линии.

\textbf{Оборудование:} осциллограф АКТАКОМ ADS-6142H; генератор АКИП~3420/1;
бухта с коаксиальным кабелем РК~50-4-11; схемный блок <<модель длинной
линии>>; магазин сопротивлений Р33; соединительные провода.

В работе по резонансным частотам, при которых сигналы в начале и в конце
кабеля совпадают по фазе, оцениваются фазовая и групповая скорости
распространения сигнала. Сняты АЧХ и ФЧХ кабеля длиной $l\approx50$~м в
диапазоне $1\div40$~МГц; по ним определены погонные индуктивность и ёмкость,
диэлектрическая и магнитная проницаемости заполнения кабеля, а по
частотной зависимости затухания (скин-эффект) "--- удельная проводимость
материала жилы. На LC-цепочке, моделирующей длинную линию, измерены
дисперсионная зависимость $\Delta\varphi(\nu)$, резонансные частоты при
короткозамкнутом и разомкнутом выходе и распределение амплитуды вдоль
цепочки.

\section{Теоретические сведения}

Все формулы ниже записаны в СИ (в методичке "--- в СГС; переход
осуществляется заменами $L_x/c^2\to L'$, $C_x\to C'$, $\varepsilon\to
4\pi\varepsilon_0\varepsilon$, $\mu\to\mu_0\mu/4\pi$). Штрихом обозначены
погонные (на единицу длины) величины.

\subsection{Погонные параметры коаксиального кабеля}
Кабель "--- центральная жила диаметра $d = 2r_1$ внутри проводящей оболочки
внутреннего диаметра $D = 2r_2$; пространство между ними заполнено средой
с проницаемостями $\varepsilon$, $\mu$. Погонные индуктивность и ёмкость:
\begin{equation}
    \label{eq:LC}
    L' = \frac{\mu_0\mu}{2\pi}\ln\frac{D}{d},
    \qquad
    C' = \frac{2\pi\varepsilon_0\varepsilon}{\ln(D/d)}.
\end{equation}
Погонное сопротивление проводников $R' = 1/(\sigma S)$, где $\sigma$ "---
удельная проводимость, $S$ "--- площадь сечения, по которому течёт ток.

\subsection{Телеграфные уравнения}
Для элемента $dx$ линии падение напряжения складывается из ЭДС
самоиндукции и омического падения, а изменение тока связано с зарядкой
погонной ёмкости:
\begin{equation}
    \label{eq:telegraph}
    \frac{\partial U}{\partial x} = -L'\frac{\partial I}{\partial t} - R'I,
    \qquad
    \frac{\partial I}{\partial x} = -C'\frac{\partial U}{\partial t}.
\end{equation}
Исключая ток, получаем волновое уравнение с затуханием:
\begin{equation}
    \label{eq:wave}
    \frac{\partial^2 U}{\partial t^2}
    - V_\text{ф}^2\frac{\partial^2 U}{\partial x^2}
    + \gamma\frac{\partial U}{\partial t} = 0,
    \qquad
    V_\text{ф} = \frac{1}{\sqrt{L'C'}} = \frac{c}{\sqrt{\varepsilon\mu}},
    \qquad
    \gamma = R'C'V_\text{ф}^2 = \frac{R'}{L'}.
\end{equation}

\subsection{Бегущая волна, затухание и волновое сопротивление}
Решение ищем в виде $U = U_0 e^{-i\omega t}e^{(ik-\alpha)x}$. Подстановка
в~\eqref{eq:wave} даёт дисперсионное соотношение
\begin{equation}
    \label{eq:dispersion}
    \omega^2 = V_\text{ф}^2\left(k^2-\alpha^2\right),
    \qquad
    2\alpha k V_\text{ф}^2 = \omega\gamma
    \;\;\Longleftrightarrow\;\;
    2\alpha k = \omega R'C'.
\end{equation}
При малом затухании ($\alpha\ll k$)
\begin{equation}
    \label{eq:alpha_k}
    k = \frac{\omega}{V_\text{ф}},
    \qquad
    \alpha = \frac{R'C'V_\text{ф}}{2} = \frac{R'}{2R_0},
\end{equation}
где $R_0$ "--- волновое сопротивление (отношение $U/I$ в бегущей волне):
\begin{equation}
    \label{eq:R0}
    R_0 = \sqrt{\frac{L'}{C'}}.
\end{equation}
Если линия нагружена на $R = R_0$, бегущая волна <<не замечает>> конца
линии и отражения нет. При $R\ne R_0$ (в частности, при $R\to\infty$ и
$R = 0$) возникает отражённая волна, и в линии устанавливается
(частично) стоячая волна.

Для согласованной линии длины $l$ амплитуда на нагрузке и набег фазы равны
\begin{equation}
    \label{eq:Un_dphi}
    U_\text{н} = U_0 e^{-\alpha l},
    \qquad
    \Delta\varphi = kl,
\end{equation}
откуда по измеренным амплитудам и сдвигу фаз
\begin{equation}
    \label{eq:alpha_k_exp}
    \alpha(\omega) = \frac{1}{l}\ln\frac{U_0}{U_\text{н}},
    \qquad
    k(\omega) = \frac{\Delta\varphi}{l}.
\end{equation}
Осциллограф позволяет измерить лишь сдвиг $\Delta x < T$ между
максимумами сигналов, т.\,е. часть фазы, меньшую $2\pi$. Поскольку
$\Delta\varphi(\omega)$ монотонно растёт, полный набег фазы
\begin{equation}
    \label{eq:phase_unwrap}
    \Delta\varphi = 2\pi\left(\nu\,\Delta x + m\right),
\end{equation}
где целое $m$ "--- число целых периодов, уложившихся во время
распространения (для $l\approx50$~м $m$ увеличивается на единицу примерно
каждые 4~МГц).

\subsection{Резонансный метод измерения скоростей}
Сигналы в начале и в конце согласованной линии совпадают по фазе, когда
набег фазы кратен $2\pi$: $kl = 2\pi n$, т.\,е.
\begin{equation}
    \label{eq:res_matched}
    \nu_n = n\,\frac{V_\text{ф}}{l}
    \quad\Longrightarrow\quad
    V_\text{ф} = l\,\Delta\nu,
\end{equation}
где $\Delta\nu = \nu_{n+1}-\nu_n$ "--- шаг между соседними резонансами.
Постоянство шага означает отсутствие дисперсии в исследуемом диапазоне.
При разомкнутом конце волна отражается и возвращается к генератору, где
поглощается его согласованным (50~Ом) выходным сопротивлением. Напряжение
на входе $U_0\bigl(1+e^{2ikl}\bigr) = 2U_0\cos(kl)\,e^{ikl}$, на выходе
$2U_0e^{ikl}$: условие совпадения экстремумов и нулей сигналов на входе и
на выходе ($kl = 2\pi n$, амплитуда на входе максимальна) остаётся
тем же, поэтому и здесь $V_\text{ф} = l\,\Delta\nu$; отражённая волна
меняет лишь амплитуду сигнала на входе.

Для последовательности прямоугольных импульсов совпадение импульсов в
начале и в конце линии наступает, когда время распространения кратно
периоду повторения: $l/V_\text{гр} = nT = n/\nu_n$, поэтому
$V_\text{гр} = l\,\Delta\nu$. Здесь измеряется скорость переноса
огибающей "--- групповая скорость $V_\text{гр} = d\omega/dk$; в отсутствие
дисперсии $V_\text{гр} = V_\text{ф}$.

\subsection{Определение $L'$, $C'$, $\varepsilon$, $\mu$}
Первое из уравнений~\eqref{eq:dispersion} перепишем как линейную
зависимость
\begin{equation}
    \label{eq:y1}
    y_1 = L'C'\,x_1,
    \qquad
    x_1 = \omega^2,
    \quad
    y_1 = k^2(\omega)-\alpha^2(\omega).
\end{equation}
Наклон прямой даёт $L'C'$, а с учётом~\eqref{eq:R0} при известном
$R_0$
\begin{equation}
    \label{eq:LC_from_slope}
    L' = R_0\sqrt{L'C'},
    \qquad
    C' = \frac{\sqrt{L'C'}}{R_0},
    \qquad
    \varepsilon\mu = c^2 L'C'.
\end{equation}
Затем из~\eqref{eq:LC} по известным $d$ и $D$ находятся $\varepsilon$ и
$\mu$ материала, заполняющего кабель.

\subsection{Скин-эффект и проводимость проводников}
Толщина скин-слоя
\begin{equation}
    \label{eq:skin}
    \delta = \frac{1}{\sqrt{\pi\nu\mu_0\sigma}}
\end{equation}
для меди уже при $\nu = 1$~МГц составляет $\approx65$~мкм, что примерно в
десять раз меньше радиуса жилы ($d = 1{,}37$~мм). Считая, что ток течёт в
приповерхностном слое жилы, и вычисляя джоулевы потери, получаем
погонное сопротивление
\begin{equation}
    \label{eq:Rx_skin}
    R' = \frac{2}{\sigma\pi d\,\delta}
       = \frac{2}{d}\sqrt{\frac{\mu_0\nu}{\pi\sigma}}.
\end{equation}

\paragraph{Метод А.} Подставляя~\eqref{eq:Rx_skin}
в~\eqref{eq:alpha_k}:
\begin{equation}
    \label{eq:methodA}
    y_2 = \alpha(\omega) = s_2\,x_2,
    \quad x_2 = \sqrt{\nu},
    \qquad
    s_2 = \frac{1}{dR_0}\sqrt{\frac{\mu_0}{\pi\sigma}}
    \;\;\Longrightarrow\;\;
    \sigma = \frac{\mu_0}{\pi\left(s_2\,d\,R_0\right)^2}.
\end{equation}

\paragraph{Метод Б.} Из второго уравнения~\eqref{eq:dispersion}
$\alpha k = \omega R'C'/2$, т.\,е.
\begin{equation}
    \label{eq:methodB}
    y_3 = \alpha(\omega)\,k(\omega) = s_3\,x_3,
    \quad x_3 = \nu^{3/2},
    \qquad
    s_3 = \frac{2\pi C'}{d}\sqrt{\frac{\mu_0}{\pi\sigma}}
    \;\;\Longrightarrow\;\;
    \sigma = \frac{\mu_0}{\pi}\left(\frac{2\pi C'}{d\,s_3}\right)^2.
\end{equation}
Метод Б не требует знания $R_0$ и использует как АЧХ, так и ФЧХ.

\subsection{Модель длинной линии}
Модель "--- цепочка из $N = 10$ звеньев с индуктивностью $L$ и ёмкостью
$C$ (\cref{fig:model_scheme}). Для волны $U_n\propto e^{i(\omega t-n\theta)}$
(где $\theta$ "--- сдвиг фаз на одно звено) из законов Кирхгофа
получается дисперсионное соотношение
\begin{equation}
    \label{eq:chain_dispersion}
    \sin\frac{\theta}{2} = \frac{\omega\sqrt{LC}}{2} = \frac{\nu}{\nu_0},
    \qquad
    \nu_0 = \frac{1}{\pi\sqrt{LC}}.
\end{equation}
При $\nu\ll\nu_0$ $\theta\approx2\pi\nu\sqrt{LC}$, что
соответствует~\eqref{eq:Un_dphi} для непрерывной линии с $L' = L/a$,
$C' = C/a$ на звено длиной $a$. С ростом частоты зависимость $\theta(\nu)$
становится нелинейной (дисперсия), а при $\nu>\nu_0$ волна по цепочке не
распространяется: $\nu_0$ "--- предельная частота. Согласованная нагрузка
при $\nu\ll\nu_0$
\begin{equation}
    \label{eq:R0_model}
    R_0 = \sqrt{\frac{L}{C}}.
\end{equation}

\paragraph{Резонансы модели.} При $R = 0$ или $R\to\infty$ волна
полностью отражается от выхода и в цепочке устанавливается стоячая волна.
Выходное сопротивление генератора (50~Ом) много меньше $R_0$, поэтому для
стоячей волны вход можно считать закороченным. Тогда условие резонанса
\begin{equation}
    \label{eq:model_res}
    R = 0:\;\; N\theta_m = m\pi,
    \qquad
    R\to\infty:\;\; N\theta_m = \left(m-\tfrac12\right)\pi,
    \qquad
    \nu_m = \nu_0\sin\frac{\theta_m}{2}.
\end{equation}
Распределение амплитуды вдоль цепочки на $m$-м резонансе имеет вид
$|U_n|\propto|\sin(n\theta_m)|$ (вход "--- узел напряжения): при
$R = 0$ узел есть и на выходе, при $R\to\infty$ на выходе пучность;
число пучностей растёт с номером резонанса.

\section{Экспериментальная установка}

\subsection{Коаксиальный кабель}
Один конец кабеля через тройник подключён к выходу генератора и к входу
CH1 осциллографа (тройник "--- непосредственно на входе осциллографа,
чтобы фазовый сдвиг между входом осциллографа и началом линии был
минимален), второй конец "--- ко входу CH2 (\cref{fig:setup_cable}).
Входное сопротивление канала CH2 (50~Ом или 1~МОм) играет роль нагрузки
линии: 50~Ом "--- согласованная нагрузка, 1~МОм "--- разомкнутый конец.

\begin{figure}[!ht]
    \centering
    \labimage{setup_cable.png}{0.7}
    \caption{Схема установки для наблюдения распространения сигналов вдоль
             коаксиального кабеля}
    \label{fig:setup_cable}
\end{figure}

\subsection{Модель длинной линии}
Цепочка из 10 катушек $L = 1{,}5$~мГн и конденсаторов $C = 47$~нФ
(крайние конденсаторы $C_1 = C/2 = 23{,}5$~нФ). Сигнал генератора подаётся
на вход, к выходу подключается магазин сопротивлений Р33; CH1 "--- вход
цепочки, CH2 "--- любая из клемм $0\ldots10$ (\cref{fig:model_scheme}).

\begin{figure}[!ht]
    \centering
    \begin{circuitikz}[scale=0.9, transform shape]
        \draw (0,0) to[sV, l=Г] (0,2.5);
        \draw (0,2.5) -- (1,2.5);
        \draw (1,2.5) to[C, l_=$C/2$] (1,0);
        \draw (1,2.5) to[L, l=$L$] (3.5,2.5);
        \draw (3.5,2.5) to[C, l_=$C$] (3.5,0);
        \draw (3.5,2.5) to[L, l=$L$] (6,2.5);
        \draw (6,2.5) to[C, l_=$C$] (6,0);
        \draw (6,2.5) -- (6.6,2.5);
        \node at (7.25,2.5) {$\cdots$};
        \node at (7.25,0) {$\cdots$};
        \draw (7.9,2.5) to[L, l=$L$] (10,2.5);
        \draw (10,2.5) to[C, l_=$C/2$] (10,0);
        \draw (10,2.5) -- (11.5,2.5) to[R, l=$R_\text{н}$] (11.5,0);
        \draw (0,0) -- (6.6,0);
        \draw (7.9,0) -- (11.5,0);
        \node[above] at (1,2.6) {0};
        \node[above] at (3.5,2.6) {1};
        \node[above] at (6,2.6) {2};
        \node[above] at (10,2.6) {10};
    \end{circuitikz}
    \caption{Модель длинной линии: 10 звеньев $L = 1{,}5$~мГн,
             $C = 47$~нФ, нагрузка "--- магазин сопротивлений $R_\text{н}$.
             Цифрами обозначены номера клемм}
    \label{fig:model_scheme}
\end{figure}


\begin{table}[!ht]
    \centering
    \caption{Параметры установки (из \filename{data/raw/setup_params.csv})}
    \label{tab:setup}
    \begin{tblr}{colspec={|l|c|c|}, hlines, vlines, row{1}={font=\bfseries}}
        Величина & Обозначение & Значение \\
        Длина кабеля & $l$ & 50~м \\
        Волновое сопротивление кабеля (номинал) & $R_0$ & 50~Ом \\
        Диаметр центральной жилы & $d$ & 1,37~мм \\
        Внутренний диаметр оплётки & $D$ & 4,6~мм \\
        Индуктивность звена модели & $L$ & 1,5~мГн \\
        Ёмкость звена модели & $C$ & 47~нФ \\
        Число звеньев & $N$ & 10 \\
        Звеньев между клеммами при измерении $\Delta\varphi$ & $K$ & 5 \\
    \end{tblr}
\end{table}

\section{Ход работы и результаты}

\subsection{Качественные наблюдения}
На генераторе устанавливался синусоидальный сигнал амплитудой 4~В
(смещение 0, фаза $0^\circ$, нагрузка 50~Ом), затем последовательность
прямоугольных импульсов (амплитуда 4~В, смещение 2~В, длительность
10~нс, фронт 4~нс). Частота менялась в пределах $1\div30$~МГц, при
согласованной (50~Ом) и разомкнутой (1~МОм) нагрузке.

При согласованной нагрузке сигнал на конце линии повторяет по форме
входной: синусоида на CH2 отстаёт от синусоиды на CH1, и с ростом частоты
отставание по фазе монотонно увеличивается, проходя через $2\pi$ примерно
каждые 4~МГц (время пробега $l/V\approx0{,}25$~мкс не зависит от частоты).
Амплитуда на конце линии несколько меньше входной и медленно убывает с
частотой "--- проявление омических потерь (скин-эффекта). Амплитуда на
входе от частоты практически не зависит, т.\,е. отражённой волны нет.

При разомкнутом конце (1~МОм) амплитуда сигнала на конце линии примерно
вдвое больше, чем при согласованной нагрузке (падающая и отражённая волны
складываются), а амплитуда на входе периодически меняется с частотой:
она максимальна, когда сигналы в начале и в конце совпадают по фазе, и
минимальна в промежуточных точках (интерференция падающей волны и волны,
вернувшейся ко входу).

Прямоугольные импульсы доходят до конца линии с задержкой
$\approx0{,}25$~мкс; их фронты немного затягиваются, а амплитуда
уменьшается из-за большего затухания высокочастотных составляющих
спектра. При разомкнутом конце на входе помимо исходного виден
отражённый импульс, вернувшийся через $2l/V\approx0{,}5$~мкс; при
согласованной нагрузке отражения нет.

\begin{figure}[!ht]
    \centering
    \begin{subfigure}{0.48\textwidth}
        \centering
        \labimage{osc_sine_matched.png}{1}
        \caption{синусоида, 50~Ом}
    \end{subfigure}\hfill
    \begin{subfigure}{0.48\textwidth}
        \centering
        \labimage{osc_sine_open.png}{1}
        \caption{синусоида, 1~МОм}
    \end{subfigure}

    \vspace{3mm}
    \begin{subfigure}{0.48\textwidth}
        \centering
        \labimage{osc_pulse_matched.png}{1}
        \caption{импульсы, 50~Ом}
    \end{subfigure}\hfill
    \begin{subfigure}{0.48\textwidth}
        \centering
        \labimage{osc_pulse_open.png}{1}
        \caption{импульсы, 1~МОм}
    \end{subfigure}
    \caption{Характерные осциллограммы сигналов в начале (CH1) и в конце
             (CH2) кабеля}
    \label{fig:osc_qual}
\end{figure}

\subsection{Фазовая и групповая скорости}
Измеренные резонансные частоты (сигналы в начале и в конце линии
совпадают по фазе) приведены в \cref{tab:cable_res}, зависимость $\nu_n(n)$
"--- на \cref{fig:cable_res}. Шаг $\Delta\nu$ найден МНК как наклон прямой
$\nu_n = \Delta\nu\cdot n + b$.

% Таблица резонансов одной серии: #1 -- файл, #2 -- подпись
\newcommand{\restable}[2]{%
    \begin{minipage}[t]{0.24\textwidth}
        \centering\small #2\\[1mm]
        \csvreader[
            tabularray = {colspec={|c|c|}, hlines, vlines,
                          row{1}={font=\bfseries}},
            table head = {$n$ & $\nu_n$, МГц \\},
        ]{../data/processed/#1}{n=\csvN, nu_MHz=\csvNu}{$\csvN$ & $\csvNu$}
    \end{minipage}%
}

\begin{table}[!ht]
    \centering
    \caption{Резонансные частоты коаксиального кабеля}
    \label{tab:cable_res}
    \restable{cable_res_sine_matched.csv}{синус, 50~Ом}\hfill
    \restable{cable_res_sine_open.csv}{синус, 1~МОм}\hfill
    \restable{cable_res_pulse_matched.csv}{импульсы, 50~Ом}\hfill
    \restable{cable_res_pulse_open.csv}{импульсы, 1~МОм}
\end{table}

% Точки серий сдвинуты по горизонтали на +-0,05..0,15, чтобы они не
% закрывали друг друга (номер резонанса -- целое число).
\begin{figure}[!ht]
    \centering
    \begin{subfigure}{\textwidth}
        \centering
        \begin{tikzpicture}
            \begin{axis}[
                labplot, height=0.45\textwidth,
                xlabel={$n$},
                ylabel={$\nu_n$, МГц},
                xmin=0, xmax=7.6, ymin=0, ymax=30,
                xtick={1,...,7},
                legend pos=north west,
            ]
                \addplot[thin, gray, no marks, domain=0:7.6, forget plot]
                    {\resnum{dnuSineMatched}*x};
                \addplot[only marks, mark=*, blue]
                    table[col sep=comma, x expr=\thisrow{n}-0.15, y=nu_MHz]
                    {../data/processed/cable_res_sine_matched.csv};
                \addlegendentry{синус, 50~Ом}
                \addplot[only marks, mark=o, red, thick]
                    table[col sep=comma, x expr=\thisrow{n}-0.05, y=nu_MHz]
                    {../data/processed/cable_res_sine_open.csv};
                \addlegendentry{синус, 1~МОм}
                \addplot[only marks, mark=square*, black]
                    table[col sep=comma, x expr=\thisrow{n}+0.05, y=nu_MHz]
                    {../data/processed/cable_res_pulse_matched.csv};
                \addlegendentry{импульсы, 50~Ом}
                \addplot[only marks, mark=triangle, green!50!black, thick]
                    table[col sep=comma, x expr=\thisrow{n}+0.15, y=nu_MHz]
                    {../data/processed/cable_res_pulse_open.csv};
                \addlegendentry{импульсы, 1~МОм}
            \end{axis}
        \end{tikzpicture}
        \caption{резонансные частоты; прямая "--- $\nu = \Delta\nu\cdot n$
                 для синуса при 50~Ом}
    \end{subfigure}

    \vspace{3mm}
    \begin{subfigure}{\textwidth}
        \centering
        \begin{tikzpicture}
            \begin{axis}[
                labplot, height=0.45\textwidth,
                xlabel={$n$},
                ylabel={$\nu_n/n$, МГц},
                xmin=0, xmax=7.6, ymin=3.85, ymax=4.1,
                xtick={1,...,7},
                ytick={3.85,3.9,...,4.1},
                yticklabel style={/pgf/number format/fixed,
                                  /pgf/number format/precision=2},
                legend style={at={(0.5,-0.28)}, anchor=north},
                legend columns=3,
            ]
                \addplot[thin, gray, dashed, no marks, domain=0:7.6]
                    {\resnum{dnuSineMatched}};
                \addlegendentry{$\Delta\nu$ (синус, 50~Ом)}
                \addplot[only marks, mark=*, blue]
                    table[col sep=comma, x expr=\thisrow{n}-0.15, y=nu_per_n]
                    {../data/processed/cable_res_sine_matched.csv};
                \addlegendentry{синус, 50~Ом}
                \addplot[only marks, mark=o, red, thick]
                    table[col sep=comma, x expr=\thisrow{n}-0.05, y=nu_per_n]
                    {../data/processed/cable_res_sine_open.csv};
                \addlegendentry{синус, 1~МОм}
                \addplot[only marks, mark=square*, black]
                    table[col sep=comma, x expr=\thisrow{n}+0.05, y=nu_per_n]
                    {../data/processed/cable_res_pulse_matched.csv};
                \addlegendentry{импульсы, 50~Ом}
                \addplot[only marks, mark=triangle, green!50!black, thick]
                    table[col sep=comma, x expr=\thisrow{n}+0.15, y=nu_per_n]
                    {../data/processed/cable_res_pulse_open.csv};
                \addlegendentry{импульсы, 1~МОм}
            \end{axis}
        \end{tikzpicture}
        \caption{частота, приходящаяся на один резонанс, $\nu_n/n = V/l$}
    \end{subfigure}
    \caption{Резонансные частоты кабеля для четырёх режимов. Точки разных
             серий сдвинуты по горизонтали на $\pm0{,}05\div0{,}15$, чтобы
             не перекрываться: в~(а) все серии ложатся на одну прямую,
             в~(б) видно, что $\nu_n/n$ во всех режимах постоянно с
             точностью $\approx1{,}5$~\%}
    \label{fig:cable_res}
\end{figure}

Результаты (\cref{eq:res_matched}, $l = 50$~м) приведены в
\cref{tab:velocities}. Частоты импульсных серий записывались с шагом
1~МГц и точно легли на прямую, поэтому их погрешность оценена по
разрешению записи частоты ($\pm0{,}05$~МГц).

\begin{table}[!ht]
    \centering
    \caption{Скорости распространения сигнала в кабеле}
    \label{tab:velocities}
    \begin{tblr}{colspec={|l|c|c|c|}, hlines, vlines, row{1}={font=\bfseries}}
        Серия & $\Delta\nu$, МГц & $V$, м/с & $V/c$ \\
        синус, 50~Ом ($V_\text{ф}$) & $\res{dnuSineMatched}$ & $\res{VSineMatched}$ & $\res{VratioSineMatched}$ \\
        синус, 1~МОм ($V_\text{ф}$) & $\res{dnuSineOpen}$ & $\res{VSineOpen}$ & $\res{VratioSineOpen}$ \\
        импульсы, 50~Ом ($V_\text{гр}$) & $\res{dnuPulseMatched}$ & $\res{VPulseMatched}$ & $\res{VratioPulseMatched}$ \\
        импульсы, 1~МОм ($V_\text{гр}$) & $\res{dnuPulseOpen}$ & $\res{VPulseOpen}$ & $\res{VratioPulseOpen}$ \\
    \end{tblr}
\end{table}

Во всех четырёх режимах резонансные частоты с хорошей точностью
кратны первой: зависимость $\nu_n(n)$ линейна, а отношение $\nu_n/n$
(\cref{fig:cable_res}\,б) не обнаруживает систематического хода с
номером резонанса "--- разброс не превышает $\approx1{,}5$~\% (наибольшее
отклонение "--- у первого резонанса, где относительная погрешность
отсчёта частоты максимальна). Следовательно,
в диапазоне $4\div28$~МГц дисперсия в кабеле отсутствует (в пределах
точности опыта). Фазовая скорость, найденная по синусоидальному сигналу,
и групповая, найденная по импульсам, совпадают в пределах погрешности:
$V_\text{ф}\approx V_\text{гр}\approx2{,}0\cdot10^8$~м/с $\approx
0{,}67\,c$, что и ожидается при отсутствии дисперсии. Это значение
совпадает со скоростью света в полиэтилене ($\varepsilon = 2{,}25$,
$c/\sqrt{\varepsilon} = 2{,}0\cdot10^8$~м/с), которым заполнен кабель
РК~50-4-11. Характер нагрузки (50~Ом или 1~МОм) на положение резонансов
не влияет, как и следует из теории.

\subsection{АЧХ и ФЧХ кабеля}
При согласованной нагрузке (50~Ом) с помощью курсоров измерялись
амплитуды $U_0$, $U_\text{н}$ сигналов в начале и в конце линии и
сдвиг $\Delta x$ между их максимумами (амплитуды "--- показания
осциллографа; в расчёт входит лишь их отношение). Целое число периодов
$m$ подбиралось так, чтобы набег фазы рос с частотой непрерывно. Полный набег фазы рассчитан по~\eqref{eq:phase_unwrap},
$\alpha$ и $k$ "--- по~\eqref{eq:alpha_k_exp}. Результаты "--- в
\cref{tab:afc}, АЧХ и ФЧХ "--- на \cref{fig:afc}.

\begin{table}[!ht]
    \centering
    \small
    \caption{АЧХ и ФЧХ коаксиального кабеля ($K = U_\text{н}/U_0$,
             $V_\text{ф} = \omega/k$)}
    \label{tab:afc}
    \csvreader[
        tabularray = {
            colspec = {|c|c|c|c|c|c|c|c|c|c|},
            hlines, vlines,
            row{1} = {font=\bfseries},
        },
        table head = {
            $\nu$, МГц & $U_0$, В & $U_\text{н}$, В & $\Delta x$, нс & $m$
            & $K$ & $\Delta\varphi$, рад & $\alpha$, м$^{-1}$
            & $k$, м$^{-1}$ & $V_\text{ф}$, $10^8$~м/с \\
        },
    ]{../data/processed/afc_table.csv}%
    {nu_MHz=\cNu, U0_V=\cUo, Un_V=\cUn, dx_ns=\cDx, m=\cM, Kamp=\cK,
     dphi_rad=\cPhi, alpha_m=\cAl, k_m=\cKw, Vph_1e8=\cV}%
    {$\cNu$ & $\cUo$ & $\cUn$ & $\cDx$ & $\cM$ & $\cK$ & $\cPhi$ & $\cAl$
     & $\cKw$ & $\cV$}
\end{table}

\begin{figure}[!ht]
    \centering
    \begin{subfigure}{0.49\textwidth}
        \centering
        \begin{tikzpicture}
            \begin{axis}[labplot, width=\linewidth, height=0.8\linewidth,
                         xlabel={$\nu$, МГц}, ylabel={$K = U_\text{н}/U_0$},
                         xmin=0, ymin=0]
                \addplot[mark=*, mark size=1.5pt]
                    table[col sep=comma, x=nu_MHz, y=Kamp]
                    {../data/processed/afc_processed.csv};
            \end{axis}
        \end{tikzpicture}
        \caption{АЧХ}
    \end{subfigure}\hfill
    \begin{subfigure}{0.49\textwidth}
        \centering
        \begin{tikzpicture}
            \begin{axis}[labplot, width=\linewidth, height=0.8\linewidth,
                         xlabel={$\nu$, МГц}, ylabel={$\Delta\varphi$, рад},
                         xmin=0, ymin=0]
                \addplot[mark=*, mark size=1.5pt]
                    table[col sep=comma, x=nu_MHz, y=dphi_rad]
                    {../data/processed/afc_processed.csv};
            \end{axis}
        \end{tikzpicture}
        \caption{ФЧХ}
    \end{subfigure}
    \caption{Амплитудно- и фазово-частотная характеристики кабеля}
    \label{fig:afc}
\end{figure}

Коэффициент передачи $K$ плавно убывает с частотой от $\approx0{,}9$ при
1~МГц до $\approx0{,}7$ при 40~МГц, т.\,е. затухание растёт с частотой,
как и должно быть при скин-эффекте ($\alpha\propto\sqrt\nu$). Ступенчатый
вид АЧХ объясняется дискретностью отсчёта амплитуд курсорами
осциллографа. ФЧХ линейна во всём диапазоне: набег фазы при 40~МГц
составляет $\approx63$~рад $\approx10\cdot2\pi$, а фазовая скорость
$V_\text{ф} = \omega/k$ во всех точках \cref{tab:afc} равна
$(1{,}97\div2{,}02)\cdot10^8$~м/с без систематической зависимости от
частоты (выпадает лишь точка 7~МГц, $2{,}09\cdot10^8$~м/с, по-видимому
из-за ошибки отсчёта $\Delta x$). Дисперсии нет и в диапазоне до 40~МГц.

\subsection{Погонные параметры кабеля}
Зависимость $y_1(x_1)$~\eqref{eq:y1} показана на \cref{fig:y1}.

\begin{figure}[!ht]
    \centering
    \begin{tikzpicture}
        \begin{axis}[
            labplot,
            xlabel={$x_1 = \omega^2$, $10^{15}$~рад$^2$/с$^2$},
            ylabel={$y_1 = k^2-\alpha^2$, м$^{-2}$},
            xmin=0, ymin=0,
            legend pos=north west,
        ]
            \addplot[only marks, mark=*]
                table[col sep=comma, x=x1, y=y1]
                {../data/processed/afc_processed.csv};
            \addlegendentry{эксперимент}
            \addplot[thin, no marks]
                table[col sep=comma, x=x1, y=y1_fit]
                {../data/processed/afc_processed.csv};
            \addlegendentry{МНК}
        \end{axis}
    \end{tikzpicture}
    \caption{Зависимость $k^2-\alpha^2$ от $\omega^2$}
    \label{fig:y1}
\end{figure}

По наклону прямой и~\eqref{eq:LC_from_slope}, \eqref{eq:LC}
($R_0 = 50$~Ом, $\ln(D/d) = 1{,}21$):
\begin{gather*}
    L'C' = \res{LCprod}~\text{с}^2/\text{м}^2,
    \\
    V_\text{ф} = \res{Vph}~\text{м/с} = \res{Vphratio}\,c,
    \\
    L' = \res{Lx}~\text{Гн/м},
    \qquad
    C' = \res{Cx}~\text{Ф/м},
    \\
    \varepsilon = \res{eps},
    \qquad
    \mu = \res{mu},
    \qquad
    \varepsilon\mu = \res{epsmu}.
\end{gather*}

Экспериментальные точки с высокой точностью ложатся на прямую,
проходящую через начало координат (свободный член
$\res{b1}$~м$^{-2}$ равен нулю в пределах погрешности). Погонная
ёмкость $C'\approx100$~пФ/м совпадает с номинальной для кабеля
РК~50-4-11 ($C'_\text{ном}\approx100$~пФ/м). Погонная индуктивность
$L'\approx2{,}51\cdot10^{-7}$~Гн/м на $\approx4$~\% больше вычисленной
по геометрии при $\mu = 1$ ($\res{LxNom}$~Гн/м); соответственно
$\varepsilon = 2{,}19$ немного меньше табличного значения для
полиэтилена ($2{,}25$), а $\mu\approx1{,}04$ немного больше единицы.
Произведение $\varepsilon\mu = 2{,}26$, определяемое непосредственно
наклоном прямой, совпадает с табличным $\varepsilon = 2{,}25$ для
полиэтилена ($\mu = 1$) с точностью 0,6~\%. Разделение $\varepsilon\mu$ на
сомножители опирается на номинальные $R_0 = 50$~Ом и геометрию кабеля
$\ln(D/d)$, известные с точностью в несколько процентов, поэтому
отклонения $\varepsilon$ и $\mu$ от табличных значений на $3\div4$~\%
следует считать несущественными: материал диэлектрика немагнитен
($\mu = 1$), $\varepsilon\approx2{,}2$. Указанные погрешности
$\varepsilon$, $\mu$, $L'$, $C'$ "--- только случайные (от МНК);
систематическая погрешность, связанная с $R_0$ и $\ln(D/d)$, порядка
$2\div5$~\%.

\subsection{Удельная проводимость материала жилы}

\begin{figure}[!ht]
    \centering
    \begin{subfigure}{0.49\textwidth}
        \centering
        \begin{tikzpicture}
            \begin{axis}[labplot, width=\linewidth, height=0.8\linewidth,
                         xlabel={$\sqrt{\nu}$, $10^3$~Гц$^{1/2}$},
                         ylabel={$\alpha$, $10^{-3}$~м$^{-1}$},
                         xmin=0, ymin=0]
                \addplot[only marks, mark=*, mark size=1.5pt]
                    table[col sep=comma, x=x2, y=y2]
                    {../data/processed/afc_processed.csv};
                \addplot[thin, no marks]
                    table[col sep=comma, x=x2, y=y2_fit]
                    {../data/processed/afc_processed.csv};
            \end{axis}
        \end{tikzpicture}
        \caption{метод А: $\alpha(\sqrt{\nu})$}
    \end{subfigure}\hfill
    \begin{subfigure}{0.49\textwidth}
        \centering
        \begin{tikzpicture}
            \begin{axis}[labplot, width=\linewidth, height=0.8\linewidth,
                         xlabel={$\nu^{3/2}$, $10^9$~Гц$^{3/2}$},
                         ylabel={$\alpha k$, $10^{-3}$~м$^{-2}$},
                         xmin=0, ymin=0]
                \addplot[only marks, mark=*, mark size=1.5pt]
                    table[col sep=comma, x=x3, y=y3]
                    {../data/processed/afc_processed.csv};
                \addplot[thin, no marks]
                    table[col sep=comma, x=x3, y=y3_fit]
                    {../data/processed/afc_processed.csv};
            \end{axis}
        \end{tikzpicture}
        \caption{метод Б: $\alpha k(\nu^{3/2})$}
    \end{subfigure}
    \caption{Определение удельной проводимости жилы}
    \label{fig:sigma}
\end{figure}

По наклонам прямых (\cref{fig:sigma}) и формулам~\eqref{eq:methodA},
\eqref{eq:methodB}:
\begin{gather*}
    s_2 = \res{sTwo}~\text{м}^{-1}\text{Гц}^{-1/2},
    \\
    \sigma_\text{А} = \res{sigmaA}~\text{См/м}
    = \res{sigmaAratio}\,\sigma_\text{Cu},
    \\
    s_3 = \res{sThree}~\text{м}^{-2}\text{Гц}^{-3/2},
    \\
    \sigma_\text{Б} = \res{sigmaB}~\text{См/м}
    = \res{sigmaBratio}\,\sigma_\text{Cu},
\end{gather*}
где $\sigma_\text{Cu} = 5{,}96\cdot10^7$~См/м "--- табличное значение для
меди. Свободный член прямой метода А: $\res{bTwo}$~м$^{-1}$.

Оба метода дают значения, совпадающие друг с другом и с табличной
проводимостью меди в пределах погрешности:
$\sigma_\text{А}\approx1{,}05\,\sigma_\text{Cu}$ ($\pm8$~\%) и
$\sigma_\text{Б}\approx1{,}03\,\sigma_\text{Cu}$ ($\pm2$~\%). Свободный
член прямой метода~А равен нулю, т.\,е. затухание действительно
пропорционально $\sqrt\nu$ "--- потери определяются скин-эффектом.
Метод~Б точнее: он использует как амплитуду, так и фазу, а величина
$\alpha k$ растёт с частотой гораздо быстрее, чем $\alpha$, так что
относительный вклад погрешности отсчёта амплитуд меньше. В методе~А
основная погрешность связана с грубым (до $\approx0{,}2$~В) отсчётом
амплитуд курсорами, из-за чего точки $\alpha(\sqrt\nu)$ имеют заметный
разброс.

То, что полученное значение не меньше табличного, несколько неожиданно:
в модели учтены потери только в центральной жиле, а потери в оплётке
и в диэлектрике должны увеличивать затухание и, значит, занижать
$\sigma$. По-видимому, это компенсируется систематическими ошибками:
неточностью номинальных $R_0$ и $d$ (в формулу метода~А они входят
квадратично), неидеальным согласованием нагрузки и отражениями на
разъёмах и тройнике.

\subsection{Модель длинной линии}

\paragraph{Предварительные оценки.} По~\eqref{eq:chain_dispersion}
и~\eqref{eq:R0_model} для $L = 1{,}5$~мГн, $C = 47$~нФ:
$\nu_0 = \res{nuZero}$~кГц, $R_0 = \res{RZeroModel}$~Ом. Согласованная
нагрузка, при которой амплитуда на всех клеммах одинакова, подобрана на
магазине: $R_\text{согл} = \res{RMatchModel}$~Ом "--- в точном согласии с
расчётом. (На схеме в описании работы один из конденсаторов подписан
<<470~нФ>>; это опечатка: распределение амплитуды вдоль цепочки не
показывает никакой особенности в этом звене, и далее для всех звеньев
принято $C = 47$~нФ.)

\paragraph{Дисперсия.} При согласованной нагрузке измерялся временной
сдвиг $\Delta t$ между сигналами на клеммах, разделённых $K = 5$
звеньями (на низких частотах $\Delta t\to K\sqrt{LC} =
\res{dtLow}$~мкс, что и наблюдается: $\Delta t\approx42$~мкс при
1~кГц). Осциллограф показывает лишь сдвиг, меньший периода $T = 1/\nu$,
поэтому, как и для кабеля, полный набег фазы
\begin{equation}
    \label{eq:model_unwrap}
    \Delta\varphi_K = 2\pi\left(\nu\,\Delta t + m\right),
    \qquad
    \Delta\varphi = \frac{\Delta\varphi_K}{K},
\end{equation}
где целое $m$ подобрано так, чтобы $\Delta\varphi(\nu)$ росла
монотонно и непрерывно. При $\nu\le19$~кГц сдвиг меньше периода
($m = 0$); при 22~кГц измеренное $\Delta t = 1{,}1$~мкс резко падает,
т.\,е. полный сдвиг превысил период $T = 45$~мкс ($m = 1$, полная
задержка $46{,}6$~мкс "--- продолжение хода $44\div47$~мкс на меньших
частотах); при 35~кГц ($T = 28{,}6$~мкс) задержка превышает уже два
периода ($m = 2$, $59$~мкс), так как групповое время вблизи $\nu_0$
быстро растёт. Сдвиг фаз на одно звено $\Delta\varphi$ сравнивается
с зависимостью~\eqref{eq:chain_dispersion} для цепочки и с линейной
зависимостью $\theta = 2\pi\nu\sqrt{LC}$ для непрерывной линии
(\cref{tab:model_phase}, \cref{fig:model_phase}).

\begin{table}[!ht]
    \centering
    \caption{Сдвиг фаз в модели длинной линии: $\Delta\varphi_K$ "--- на
             $K = 5$ звеньев, $\Delta\varphi$ "--- на одно звено}
    \label{tab:model_phase}
    \csvreader[
        tabularray = {colspec={|c|c|c|c|c|c|c|}, hlines, vlines,
                      row{1}={font=\bfseries}},
        table head = {$\nu$, кГц & $\Delta t$, мкс & $m$
                      & $\Delta\varphi_K$, рад & $\Delta\varphi$, рад
                      & $\theta_\text{цеп}$, рад & $\theta_\text{непр}$, рад \\},
    ]{../data/processed/model_phase_table.csv}%
    {nu_kHz=\cNu, dt_us=\cDt, m=\cM, dphi_tot=\cTot, dphi_rad=\cPhi,
     dphi_chain=\cTh, dphi_cont=\cTc}%
    {$\cNu$ & $\cDt$ & $\cM$ & $\cTot$ & $\cPhi$ & $\cTh$ & $\cTc$}
\end{table}

\begin{figure}[!ht]
    \centering
    \begin{tikzpicture}
        \begin{axis}[
            labplot,
            xlabel={$\nu$, кГц},
            ylabel={$\Delta\varphi$, рад},
            xmin=0, xmax=40, ymin=0, ymax=3.5,
            legend pos=north west,
        ]
            \addplot[only marks, mark=*]
                table[col sep=comma, x=nu_kHz, y=dphi_rad]
                {../data/processed/model_phase_processed.csv};
            \addlegendentry{эксперимент}
            \addplot[thick, no marks]
                table[col sep=comma, x=nu_kHz, y=dphi_chain]
                {../data/processed/model_phase_theory.csv};
            \addlegendentry{$2\arcsin(\nu/\nu_0)$}
            \addplot[thin, dashed, no marks]
                table[col sep=comma, x=nu_kHz, y=dphi_cont]
                {../data/processed/model_phase_theory.csv};
            \addlegendentry{$2\pi\nu\sqrt{LC}$}
        \end{axis}
    \end{tikzpicture}
    \caption{Сдвиг фаз на одно звено модели в зависимости от частоты
             (с учётом целого числа периодов $m$)}
    \label{fig:model_phase}
\end{figure}

После учёта целого числа периодов экспериментальные точки ложатся на
единую монотонную кривую. До $\nu\approx10$~кГц она практически линейна и
совпадает с обеими теоретическими зависимостями (при $\nu\ll\nu_0$
цепочка ведёт себя как непрерывная линия). С приближением к
$\nu_0 = 37{,}9$~кГц экспериментальная зависимость, как и
теоретическая $2\arcsin(\nu/\nu_0)$, загибается вверх и заметно уходит
от прямой $2\pi\nu\sqrt{LC}$ "--- в цепочке есть дисперсия: фазовая
скорость (в звеньях за секунду) $\omega/\Delta\varphi$ убывает с ростом
частоты. Измеренные значения систематически выше теоретических на
$5\div12$~\%; наибольшее расхождение "--- вблизи $\nu_0$, где
$\Delta\varphi(\nu)$ меняется наиболее резко, а сигнал на дальних
клеммах сильно ослаблен и искажён, так что отсчёт $\Delta t$ наименее
точен. Кроме того, на расхождение влияют разброс номиналов $L$ и $C$ и
активное сопротивление катушек.

\paragraph{Резонансы.} Резонансные частоты в диапазоне от 2~кГц до
$\nu_0$ при замкнутом и разомкнутом выходе приведены в
\cref{tab:model_res} вместе с оценкой~\eqref{eq:model_res}.

% #1 -- файл, #2 -- подпись
\newcommand{\modelrestable}[2]{%
    \begin{minipage}[t]{0.48\textwidth}
        \centering\small #2\\[1mm]
        \csvreader[
            tabularray = {colspec={|c|c|c|c|}, hlines, vlines,
                          row{1}={font=\bfseries}},
            table head = {$m$ & $\nu^\text{эксп}$, кГц & $\nu^\text{теор}$, кГц
                          & $\varepsilon$, \% \\},
        ]{../data/processed/#1}%
        {n=\cN, nu_exp_kHz=\cE, nu_theor_kHz=\cT, dev_percent=\cD}%
        {$\cN$ & $\cE$ & $\cT$ & $\cD$}
    \end{minipage}%
}

\begin{table}[!ht]
    \centering
    \caption{Резонансные частоты модели длинной линии}
    \label{tab:model_res}
    \modelrestable{model_res_short.csv}{$R = 0$}\hfill
    \modelrestable{model_res_open.csv}{$R\to\infty$}
\end{table}

Порядок следования резонансов согласуется с теорией: при разомкнутом
выходе первый резонанс (четверть волны на длине цепочки) лежит ниже, чем
при замкнутом (полволны), а резонансы двух режимов чередуются:
$2{,}1 < 4{,}9 < 7{,}3 < 15{,}0 < 16{,}8$~кГц. Количественно согласие
хуже: отклонения от оценки~\eqref{eq:model_res} составляют $15\div30$~\%.
Оценка предполагает, что вход цепочки является идеальным узлом
напряжения; на самом деле выходное сопротивление генератора (50~Ом)
сравнимо с $R_0 = 179$~Ом, и входной конец лишь частично отражает
волну. Из-за этого, а также из-за потерь в катушках резонансы широкие и
слабо выраженные (особенно на высоких частотах), и положение максимума
определяется неуверенно. Поэтому удалось надёжно зарегистрировать лишь
несколько наименьших резонансов в каждом режиме; второй резонанс при
$R = 0$ ($15{,}0$~кГц) лежит между теоретическими значениями для
$m = 2$ и $m = 3$ ($11{,}7$ и $17{,}2$~кГц), т.\,е. резонанс вблизи
11--12~кГц, вероятно, был пропущен.

\paragraph{Распределение амплитуды вдоль цепочки.} Для двух наименьших
резонансов "--- $\nu_1 = \res{nuDistOne}$~кГц (выход замкнут, $R = 0$) и
$\nu_2 = \res{nuDistTwo}$~кГц (выход разомкнут, $R\to\infty$) "---
измерена амплитуда на каждой клемме (\cref{fig:model_dist}). Линиями
показаны теоретические распределения стоячей волны (узел на входе,
нормировка на максимум серии): $|\sin(n\pi/10)|$ при $R = 0$ и
$|\sin(n\pi/20)|$ при $R\to\infty$.

\begin{figure}[!ht]
    \centering
    \begin{tikzpicture}
        \begin{axis}[
            labplot,
            xlabel={номер клеммы $n$},
            ylabel={$U$, мВ},
            xmin=0, xmax=10, ymin=0, ymax=42,
            xtick={0,...,10},
            legend pos=north west,
            legend columns=2,
        ]
            \addplot[only marks, mark=*, blue] table[col sep=comma, x=cell, y=U1_mV]
                {../data/processed/model_distribution_processed.csv};
            \addlegendentry{{$\nu_1$, $R = 0$}}
            \addplot[only marks, mark=square*, red] table[col sep=comma, x=cell, y=U2_mV]
                {../data/processed/model_distribution_processed.csv};
            \addlegendentry{{$\nu_2$, $R\to\infty$}}
            \addplot[blue, thin, smooth, samples=101, domain=0:10]
                {29.8*abs(sin(deg(x*pi/10)))};
            \addlegendentry{$|\sin(n\pi/10)|$}
            \addplot[red, thin, dashed, smooth, samples=101, domain=0:10]
                {37.6*abs(sin(deg(x*pi/20)))};
            \addlegendentry{$|\sin(n\pi/20)|$}
        \end{axis}
    \end{tikzpicture}
    \caption{Распределение амплитуды напряжения вдоль цепочки на двух
             наименьших резонансных частотах}
    \label{fig:model_dist}
\end{figure}

В обоих случаях в цепочке устанавливается стоячая волна с узлом на
входе (низкоомный генератор). При замкнутом выходе ($\nu_1$) на длине
цепочки укладывается половина волны: амплитуда нарастает от входа,
достигает пучности в середине цепочки (клеммы 5--6) и снова падает
почти до нуля на закороченном выходе (остаточные 5~мВ "--- падение
напряжения на сопротивлении перемычки). При разомкнутом выходе ($\nu_2$)
укладывается четверть волны: амплитуда монотонно растёт и выходит на
пучность на открытом конце (клеммы 9--10). Форма распределений качественно
согласуется с теоретическими синусоидами; небольшой сдвиг пучности
к выходу при $R = 0$ и более пологий рост при $R\to\infty$ объясняются
тем, что вход не является идеальным узлом (см. обсуждение резонансов).

\section{Обсуждение результатов}

Скорость распространения сигнала в коаксиальном кабеле определена тремя
независимыми способами: по резонансам синусоидального сигнала
(фазовая скорость), по резонансам последовательности импульсов
(групповая скорость) и по наклону ФЧХ. Все три дают
$V\approx2{,}0\cdot10^8$~м/с $\approx0{,}67\,c$ с расхождением менее
0,5~\%, а $\nu_n/n$ и $\omega/k$ не зависят от частоты в диапазоне
$1\div40$~МГц. Это означает отсутствие дисперсии в кабеле: при
$\alpha\ll k$ (у нас $\alpha/k < 0{,}06$) дисперсионное
соотношение~\eqref{eq:dispersion} сводится к $k = \omega/V_\text{ф}$.

Погонные параметры кабеля, найденные из
зависимости $k^2-\alpha^2$ от $\omega^2$, согласуются с номинальными:
$C'\approx100$~пФ/м, $\varepsilon\mu\approx2{,}26$, что соответствует
полиэтиленовой изоляции. Основной источник систематической
погрешности здесь "--- использование номинальных $R_0$ и размеров
кабеля.

Затухание сигнала растёт с частотой как $\sqrt\nu$, что подтверждает
определяющую роль скин-эффекта в потерях. Удельная проводимость
материала жилы, найденная двумя методами, совпадает с табличной для
меди; метод~Б, использующий и АЧХ, и ФЧХ, даёт меньшую погрешность
($\approx2$~\% против $\approx8$~\%). Точность метода~А ограничена
грубым отсчётом амплитуд, которые при длине кабеля 50~м меняются лишь
на $10\div30$~\%.

На модели длинной линии подтверждены основные свойства цепочки
$LC$-звеньев: согласованная нагрузка $R_\text{согл} = \sqrt{L/C}$,
линейная зависимость сдвига фаз от частоты при $\nu\ll\nu_0$ и
сильная дисперсия вблизи предельной частоты $\nu_0$. При измерении
сдвига фаз существенно учитывать прохождение сдвига через целое число
периодов: без этого зависимость $\Delta\varphi(\nu)$ имеет ложные
скачки вниз. При полном отражении на выходе в цепочке образуются стоячие
волны, распределение амплитуды которых соответствует укладыванию
половины (выход замкнут) или четверти (выход разомкнут) волны на длине
цепочки. Резонансные частоты согласуются с простой оценкой лишь
качественно (с точностью $15\div30$~\%) из-за неидеальности граничного
условия на входе и потерь в катушках.

\section{Выводы}
\begin{enumerate}
    \item Резонансным методом измерены фазовая скорость
          $V_\text{ф} = \res{VSineMatched}$~м/с и групповая скорость
          $V_\text{гр} = \res{VPulseMatched}$~м/с распространения сигнала в
          кабеле; они совпадают и равны $\approx0{,}67\,c =
          c/\sqrt{\varepsilon_\text{пэ}}$. Дисперсия в диапазоне
          $1\div40$~МГц не обнаружена.
    \item По АЧХ и ФЧХ кабеля определены погонные параметры
          $L' = \res{Lx}$~Гн/м, $C' = \res{Cx}$~Ф/м и проницаемости
          заполнения $\varepsilon = \res{eps}$, $\mu = \res{mu}$
          ($\varepsilon\mu = \res{epsmu}$), что соответствует
          полиэтилену ($\varepsilon = 2{,}25$, $\mu = 1$) с точностью
          до систематической погрешности номинальных параметров кабеля.
    \item По частотной зависимости затухания определена удельная
          проводимость жилы: $\sigma_\text{А} = \res{sigmaA}$~См/м,
          $\sigma_\text{Б} = \res{sigmaB}$~См/м; оба значения совпадают
          в пределах погрешности с табличным для меди
          ($5{,}96\cdot10^7$~См/м).
    \item Для модели длинной линии предельная частота
          $\nu_0 = \res{nuZero}$~кГц, согласованная нагрузка
          $R_\text{согл} = \res{RMatchModel}$~Ом
          (теория $\res{RZeroModel}$~Ом). Сдвиг фаз на звено следует
          зависимости $2\arcsin(\nu/\nu_0)$ (с превышением на
          $5\div12$~\%), что демонстрирует дисперсию цепочки вблизи
          $\nu_0$. Наблюдены резонансы при замкнутом и разомкнутом выходе;
          распределения амплитуды на двух наименьших резонансах
          соответствуют стоячим волнам с половиной и четвертью длины
          волны на длине цепочки.
\end{enumerate}

\end{document}
